% TOP_PROBLEM: {"id": 30003057, "title": "Hemisphere Partition Functions and Stability Across GLSM Phases", "category_id": 16, "category": {"id": 16, "name": "physics", "display_name": "Mathematical Physics", "description": "Problems at the intersection of mathematics and physics.", "slug": "physics", "order_index": 16, "created_at": "2026-07-31T15:26:25.671Z"}, "status": "open"}
% FIRST_POSTED: 2026-10-06
% ATTEMPT_STATUS: unresolved
% OPENAI_CLAIM: {"schema_version":1,"source_repo":"https://github.com/openai/math","source_commit":"adc7f1241b42e322a6451854ab7e4b4c146bf78a","release_date":"2026-10-06","match":"related","resolution":"partial","independently_reviewed":false,"second_pass":"n/a","families":[{"family":"055","title":"Gepner symmetry and large-volume stability on threefolds","subject":"Algebraic and complex geometry","manuscripts":[{"title":"A Gepner stability condition on every smooth quintic threefold","pdf_path":"preprints/A-Gepner-stability-condition-on-every-smooth-quintic-threefold-September-24-2026/paper.pdf","readme_path":"preprints/A-Gepner-stability-condition-on-every-smooth-quintic-threefold-September-24-2026/README.md","date":"2026-09-24","theorem_ref":"thm:main (sections/introduction.tex:52)","statement_sources":[{"path":"preprints/A-Gepner-stability-condition-on-every-smooth-quintic-threefold-September-24-2026/build/sections/introduction.tex","line":52,"label":"thm:main"}]}],"lean":null,"reasoning_trace":null}]}
\documentclass[11pt]{article}
\usepackage[margin=25mm]{geometry}
\usepackage{fontspec}
\setmainfont{Latin Modern Roman}
\usepackage{amsmath,amssymb,mathtools}
\usepackage{unicode-math}
\setmathfont{Latin Modern Math}
\usepackage{enumitem}
\usepackage{xurl,hyperref,bookmark}
\hypersetup{colorlinks=true,urlcolor=blue}
\setcounter{secnumdepth}{0}
\setlength{\parindent}{0pt}
\setlength{\parskip}{5pt}
\setlength{\emergencystretch}{3em}
\begin{document}
\section*{Hemisphere Partition Functions and Stability Across GLSM Phases}

\subsection{OpenAI's claim}
Every smooth complex quintic threefold is claimed to carry a numerical Bridgeland stability condition with the normalized Gepner charge and phase shift $2/5$ under the specified autoequivalence. The cited manuscript argues for that result; its relation to this problem is specified below.

\subsection{Mathematical statement of the claimed result}
\paragraph{Mathematical claim 1.} For every smooth quintic hypersurface $X\subset\mathbb P^4_{\mathbb C}$, define on $D^b(\operatorname{Coh}X)$ the autoequivalence
\[
\Phi=T_{\mathcal O_X}\circ(-\otimes\mathcal O_X(1)),\qquad
T_{\mathcal O_X}(E)=\operatorname{Cone}\bigl(R\operatorname{Hom}(\mathcal O_X,E)\otimes\mathcal O_X\to E\bigr).
\]
OpenAI claims a numerical Bridgeland stability condition $(Z,\mathcal P)$ with support property on the full real numerical Grothendieck group such that
\[
Z(\Phi E)=e^{2\pi i/5}Z(E),\quad
\Phi(\mathcal P(\varphi))=\mathcal P(\varphi+2/5),\quad
Z(\mathcal O_x)=-1
\]
for every object $E$, real phase $\varphi$ and closed point $x$. Support means a norm and constant $C$ bound the numerical class of every semistable object by $C|Z(E)|$.

\subsection{Outline of the argument}
Manuscript 1 (family 055, 2026-09-24): Two independent perturbations of the second slope tilt control rank and first Chern character before semistability is defined. The charge bounds the remaining numerical coordinates, enabling finite refinements and Harder--Narasimhan filtrations even for a nondiscrete charge image; compatibility of the tilted aisles supplies the Gepner rotation.

\subsection{Scope relative to this problem}
The theorem produces a stability condition with the quintic Gepner symmetry and normalized eigencharge. This page instead asks that the specific hemisphere partition function define a stability condition and that the grade-restriction categories give the specified phase equivalences. The manuscript supplies neither that partition-function identification nor those commuting equivalences. It is related stability information, rather than a claimed partial solution of these GLSM assertions.

\subsection{Formal verification}
Family 055 has no Lean scope page in this snapshot; no formalization of these manuscript results is listed.

This site has not executed the Lean files or checked their dependencies. These descriptions report the scope of the linked declarations, rather than a new verification of the proofs.

\subsection{Status caveat}
OpenAI attributes these results to an internal model. Its release notice cautions that results lacking full formalization can contain mistakes. This site has not independently assessed the proofs; a solved label records the scope of the reported claim.

\subsection{Sources}
\begin{itemize}
\item OpenAI, \emph{A Gepner stability condition on every smooth quintic threefold}, 2026-09-24; citation key \texttt{OAI:A-Gepner-stability-condition-on-every-smooth-quintic-threefold-September-24-2026}. \href{https://github.com/openai/math/blob/adc7f1241b42e322a6451854ab7e4b4c146bf78a/preprints/A-Gepner-stability-condition-on-every-smooth-quintic-threefold-September-24-2026/paper.pdf}{PDF (GitHub)}; \href{https://github.com/openai/math/raw/adc7f1241b42e322a6451854ab7e4b4c146bf78a/preprints/A-Gepner-stability-condition-on-every-smooth-quintic-threefold-September-24-2026/paper.pdf}{PDF (direct)}; \href{https://github.com/openai/math/blob/adc7f1241b42e322a6451854ab7e4b4c146bf78a/preprints/A-Gepner-stability-condition-on-every-smooth-quintic-threefold-September-24-2026/README.md}{Manuscript page and citation}; \href{https://github.com/openai/math/blob/main/preprints/A-Gepner-stability-condition-on-every-smooth-quintic-threefold-September-24-2026/paper.pdf}{Latest version}; Statement source: \href{https://github.com/openai/math/blob/adc7f1241b42e322a6451854ab7e4b4c146bf78a/preprints/A-Gepner-stability-condition-on-every-smooth-quintic-threefold-September-24-2026/build/sections/introduction.tex\#L52}{thm:main}.
\item \href{https://github.com/openai/math/blob/adc7f1241b42e322a6451854ab7e4b4c146bf78a/CONTENTS.md}{OpenAI catalogue entry}.
\end{itemize}
\end{document}
